\documentclass{article}
\usepackage{sty/Layout}
\usepackage{sty/Commands}
\usepackage{sty/Titlepage}
\usepackage{siunitx}
\sisetup{separate-uncertainty = true}
\usepackage{multirow}
\usepackage{marginnote}
\usepackage{xcolor}

% ---- Link styling -------------------------------------------------------
% No red boxes: coloured link text, and in-text cross-references underlined
% so they are easy to spot on paper.
\definecolor{linkblue}{rgb}{0.10,0.35,0.65}
\hypersetup{
	colorlinks = true,
	linkcolor  = linkblue,   % ToC, \ref, \eqref, section/equation links
	citecolor  = linkblue,
	urlcolor   = linkblue,
	filecolor  = linkblue,
	linktoc    = all,
}
\usepackage[most]{tcolorbox}
\usepackage[normalem]{ulem}
\let\OrigRef\ref
\renewcommand{\ref}[1]{{\color{linkblue}\uline{\OrigRef{#1}}}}
\renewcommand{\eqref}[1]{{\color{linkblue}\uline{\textup{(\OrigRef{#1})}}}}
% -----------------------------------------------------------------------


\begin{document}

\makecoverpage{
	experiment      = 73,
	title           = {EXPERIMENT TITLE},
	group           = 8,
	authors         = {Matthes Neumann},
	tutor           = {Natalia Sychva},
	dateexperiment  = {30. August 2026}
}

\pagestyle{fancy}

\tableofcontents
\clearpage

\section{Introduction}

\subsection{Motivation}
\label{sec:motivation}

\vspace{1em}
\begin{tcolorbox}[
	colback  = linkblue!3,
	colframe = linkblue!55,
	boxrule  = 0.4pt,
	arc      = 2pt,
	left     = 6pt, right = 6pt, top = 5pt, bottom = 5pt,
	fontupper = \small,
	title    = {\scriptsize\textsf{\textbf{Nutzung von Generativer künstlicher Intelligenz}}},
	coltitle = linkblue!80!black,
	colbacktitle = linkblue!10,
	fonttitle = \normalfont,
]
	\textbf{Werkzeuge:} Claude Sonnet 5, Opus 5 (Anthropic) \\
	\textbf{Abschnitte:} Formulierungshilfen, Textgeneration und Formatierungshilfen wurden teilweise in 
	den Abschnitten \ref{sec:motivation}, \ref{sec:grundlagen}, \ref{sec:konstanten}, \ref{sec:aufbau} genutzt.
	In den restlichen Abschnitten wurden die genannten Transformatoren nur für die Recherche zu technischen
	Fragen über die Erstellung von Darstellungen in Latex und Python Bibliotheken genutzt.
\end{tcolorbox}


\subsection{Theoretische Grundlagen}
\label{sec:grundlagen}

\subsection{Versuchsaufbau}
\label{sec:aufbau}

\subsection{Verwendete Konstanten}
\label{sec:konstanten}

\section{Messprotokoll}
\label{sec:messprotokoll}


\section{Auswertung}
\label{sec:auswertung}


\section{Disskussion}
\label{sec:diskussion}


\section{Anhang}
\label{sec:diskussion}


\end{document}
